% Generated draft: requires mathematical and source review.
\documentclass[11pt]{article}
\usepackage[margin=25mm]{geometry}
\usepackage{fontspec}
\setmainfont{Latin Modern Roman}
\usepackage{amsmath,amssymb,mathtools,unicode-math}
\setmathfont{Latin Modern Math}
\usepackage{enumitem,needspace,xurl,hyperref}
\setlength{\parindent}{0pt}
\setlength{\parskip}{5pt}
\setlength{\emergencystretch}{3em}
\newcommand{\N}{\mathbb{N}}
\newcommand{\Z}{\mathbb{Z}}
\newcommand{\Q}{\mathbb{Q}}
\newcommand{\R}{\mathbb{R}}
\newcommand{\C}{\mathbb{C}}
\newcommand{\F}{\mathbb{F}}
\newcommand{\A}{\mathbb{A}}
\newcommand{\PP}{\mathbb{P}}
\newcommand{\E}{\mathbb{E}}
\newcommand{\eps}{\varepsilon}
\newcommand{\dd}{\,\mathrm{d}}
\begin{document}
\section*{UnsolvedMath ID: 3900009}
\textbf{Source problem code:} catalogue entry
\subsection{Definitions and mathematical statement}
Let $\Omega\subset\mathbb R^2$ be a bounded polygonal room and $s\in\Omega$ a point source. A point $q\in\Omega$ is illuminated from $s$ if it is reached by a finite polygonal light path starting at $s$, lying inside the room between reflections, and obeying equal angles of incidence and reflection at each side it hits. Direct paths with no reflection are allowed. Under the usual billiard convention, paths through vertices are excluded because no unique reflection rule is defined there. Let $I(\Omega,s)$ be the set of illuminated points and put
\[
U(\Omega,s)=\Omega\setminus I(\Omega,s).
\]
Characterize the pairs $(\Omega,s)$ for which $U(\Omega,s)=\varnothing$. In particular, can $U(\Omega,s)$ be nonempty for a polygonal mirrored room and a chosen source?

The source also phrases the last question as the existence of a ``dark region.'' If this means a nonempty open subset of $U(\Omega,s)$, or positive-area $U(\Omega,s)$, it is stronger than the existence of a single unilluminated point. The supplied summary attributes to Tokarsky (1995) rational-angle polygonal examples containing points not illuminated from a chosen source. That report gives a negative answer to universal illumination but does not itself supply a general characterization or establish an open dark region. Boundary targets, holes in the room, and corner trajectories are not specified in the source; the definitions above use interior targets and the usual nonvertex reflection rule.
\subsection{Short English statement}
Determine which polygonal mirrored rooms are fully illuminated from a specified point source. Distinguish the existence of an unilluminated point from the stronger possibility of an open or positive-area dark region.
\subsection{Sources}
\begin{itemize}
\item David Eppstein, \emph{The Geometry Junkyard: Open Problems}, item ``Mirrored-room illumination,'' \url{https://ics.uci.edu/~eppstein/junkyard/open.html}; accessed 2026-07-29, HTML extraction.
\item Tokarsky (1995), polygonal nonillumination examples, as attributed in the supplied research summary; no publication title or further identifier is supplied.
\item Local registry \texttt{data/unsolvedmath/problems.json}, permanent UnsolvedMath ID $3900009$, source code catalogue entry, \emph{Mirrored-room illumination}. The proposer field is ``David Eppstein and cited proposers.'' Dataset attribution: Ulam AI and the UnsolvedMath Contributors.
\end{itemize}
\end{document}
